\chapter{Operational Rigidity Principle}

\section{Purpose of Operational Rigidity}

Operational rigidity is the principle that a bounded observer cannot freely
convert local information into global structure. A refinement process may
inspect local data, update labels, split partitions, or remove possibilities,
but if each step has bounded information capacity, then global collapse
requires depth.

The principle is operational because it concerns what an admissible process can
do, not merely what is true about the object in the abstract. A global
certificate may exist, but a local refinement process may still be unable to
extract it in one step. The framework therefore separates three questions:

\begin{enumerate}
\item What global structure exists?
\item What local information is visible to the admissible process?
\item How many refinement steps are required to transfer local visibility into
global resolution?
\end{enumerate}

This chapter formulates that separation.

\section{Local Observables}

Let \(X\) be a state space. A local observable is a measurement that sees only
a bounded part of a configuration.

\begin{definition}[Local observable]
A local observable on \(X\) is a map
\[
\ell:X\to A
\]
whose value is determined by bounded local data. The codomain \(A\) is the set
of possible local types.
\end{definition}

In a graph model, \(\ell(x)\) may record the isomorphism type of a bounded
neighborhood. In a logical model, \(\ell(x)\) may record the truth values of a
bounded family of formulas. In a physical model, \(\ell(x)\) may record a
bounded detector reading.

\begin{definition}[Local type partition]
Given a local observable \(\ell:X\to A\), the associated local type partition
is
\[
x\sim_{\ell} y
\quad\Longleftrightarrow\quad
\ell(x)=\ell(y).
\]
The equivalence classes of \(\sim_{\ell}\) are the states indistinguishable by
that local observable.
\end{definition}

A local observable is useful only after its admissibility is declared. The same
mathematical object may be local in one model and nonlocal in another.

\section{Local Indistinguishability}

Local indistinguishability is the obstruction behind operational rigidity.

\begin{definition}[\(R\)-local indistinguishability]
Two configurations \(x,y\in X\) are \(R\)-locally indistinguishable if every
admissible observation of radius at most \(R\) gives the same local type on
\(x\) and \(y\).
\end{definition}

The point is not that \(x=y\). The point is that an admissible local process
cannot separate them without accumulating additional information over time.

\begin{definition}[Persistent local indistinguishability]
A family \(S_n\subseteq X_n\) has persistent local indistinguishability if,
for a fixed locality radius \(R\), many distinct configurations in \(S_n\)
share the same admissible \(R\)-local data.
\end{definition}

Persistent local indistinguishability supplies entropy. If many global
configurations have the same local view, then local observation leaves a large
residual uncertainty set.

\section{Rigidity Principle}

The operational rigidity principle is the following schematic implication:
\[
\text{large local equivalence class}
+
\text{bounded step capacity}
\Longrightarrow
\text{positive refinement depth}.
\]

\begin{theorem}[Operational Rigidity Principle]
Let \(S\subseteq X\) be a set of configurations that are indistinguishable by
the admissible local observables at the initial resolution. Suppose an
admissible refinement step can remove at most \(C>0\) units of entropy from
\(S\). If the unresolved entropy of \(S\) is \(H(S)\), then every admissible
refinement process that isolates one terminal configuration from \(S\) has
length at least
\[
\frac{H(S)}{C}.
\]
\end{theorem}

\begin{proof}
At the initial resolution, the process must distinguish among configurations
inside \(S\). The residual entropy is \(H(S)\). Since each admissible step
removes at most \(C\) entropy, after \(t\) steps the total removed entropy is
at most \(tC\). Isolation of one terminal configuration requires removal of
all \(H(S)\) units. Hence \(tC\ge H(S)\), so \(t\ge H(S)/C\).
\end{proof}

This theorem is only a bookkeeping theorem until a domain supplies a concrete
local observable, entropy functional, and capacity estimate.

\section{Local--Global Transfer}

Local--global transfer is the mechanism that turns bounded local data into a
global lower bound.

\begin{definition}[Transfer datum]
A transfer datum consists of
\[
(X,\mathcal L,\mathcal R,H,C,T),
\]
where \(X\) is a state space, \(\mathcal L\) is a class of local observables,
\(\mathcal R\) is a class of admissible refinements, \(H\) is an entropy
functional, \(C\) is a per-step capacity bound, and \(T\) is a terminal set.
\end{definition}

\begin{theorem}[Local--Global Rigidity Transfer]
Assume a transfer datum \((X,\mathcal L,\mathcal R,H,C,T)\). Let \(x_0\in X\).
Suppose the local data available to \(\mathcal L\) leaves an unresolved class
\(S(x_0)\) with entropy
\[
H(S(x_0))\ge h.
\]
Suppose every refinement in \(\mathcal R\) removes at most \(C\) units of this
unresolved entropy. Then every admissible refinement trajectory from \(x_0\)
to \(T\) has length at least
\[
\frac{h}{C}.
\]
\end{theorem}

\begin{proof}
The unresolved class \(S(x_0)\) contains the global alternatives that remain
compatible with the local data. By hypothesis, resolving the trajectory to the
terminal set requires eliminating at least \(h\) units of unresolved entropy.
The capacity bound allows at most \(C\) units to be removed per admissible
step. Therefore a trajectory of length \(t\) removes at most \(tC\) units.
Terminal resolution requires \(tC\ge h\), proving \(t\ge h/C\).
\end{proof}

\section{Graph Example}

Let \(G=(V,E)\) be a bounded-degree graph. A radius-\(R\) local observable at a
vertex \(v\) records the rooted radius-\(R\) neighborhood of \(v\). A refinement
process may update a label at \(v\) from the multiset of labels in that
neighborhood.

\begin{example}[Bounded-degree local refinement]
Assume the degree bound \(\Delta\) and radius \(R\) are fixed. Then the number
of possible rooted radius-\(R\) neighborhoods is bounded independently of the
total number of vertices. Therefore one local observation has bounded type
complexity. If a family of graphs has globally many configurations sharing
the same local type data, local refinement cannot distinguish them in one
step.
\end{example}

This example does not prove graph isomorphism lower bounds. It only explains
how local type indistinguishability becomes a URF transfer datum.

\section{Spectral Example}

In a spectral setting, local observables may be replaced by projections onto
accessible modes. The unresolved entropy is then carried by components not
separated by the available spectral information.

\begin{example}[Spectral obstruction]
Let \(L\) be a nonnegative self-adjoint operator and suppose the admissible
process only accesses a bounded spectral window. If many orthogonal directions
remain unresolved inside that window, then the refinement process must either
increase accessible capacity or take multiple steps to separate them.
\end{example}

A spectral gap supplies coercivity, but the operational question is separate:
which modes can the process actually access, and how much information can be
extracted per step?

\section{Algorithmic Example}

In an algorithmic setting, the refinement state may be a partition of a
configuration space \(\Omega\). A refinement step replaces a partition \(P_t\)
by a finer partition \(P_{t+1}\) using admissible local data.

\begin{example}[Partition refinement]
Let \(P_t\) be a partition of \(\Omega\). The unresolved entropy at time \(t\)
may be measured by the logarithm of the size of the block containing the true
configuration. If each admissible refinement splits this block into at most
\(B\) effectively distinguishable subblocks, then the entropy decrease per
step is at most \(\log B\).
\end{example}

This example is the operational form of the capacity-depth argument. It does
not assert that every algorithm is a partition refinement algorithm unless a
normalization theorem has been supplied.

\section{Failure Modes}

Operational rigidity can fail in several precise ways.

\begin{enumerate}
\item The process may be nonlocal.
\item The per-step capacity may grow with system size.
\item The local equivalence class may have small entropy.
\item The terminal condition may not require full resolution.
\item The claimed local observable may secretly encode global data.
\end{enumerate}

Each failure mode is an audit point. A valid URF argument must state why the
failure mode is absent or must record it as an open boundary.

\section{Operational Checklist}

A complete operational-rigidity application should specify:

\begin{center}
\begin{tabular}{ll}
\toprule
Field & Required content \\
\midrule
Local data & admissible observables and locality radius \\
State class & configurations under study \\
Unresolved class & configurations not separated by local data \\
Entropy & size or uncertainty of the unresolved class \\
Capacity & maximum information removed per step \\
Terminal condition & what counts as resolved \\
Transfer theorem & local data converted into global lower bound \\
\bottomrule
\end{tabular}
\end{center}

\section{Boundary}

This chapter proves only the schematic operational transfer from local
indistinguishability and bounded capacity to refinement depth. It does not
prove that every algorithm is local, does not prove unconditional graph or
complexity lower bounds, and does not prove unrestricted Chronos-RR,
unrestricted H4.1/FGL, P vs NP, or any Clay problem.
